\cotefile{77}
\scanpage{50}
\deleted{Donc si $\Delta'$ est la \unsure{sgn-section|2-sect°} de $\check{\mathbb V}(L)$ \unsure{défini|définie} dans le th, on a $x\in\Gamma\Delta$ ssi $2x\in\Gamma\Delta'$, i.e. $\Delta'=2\Delta$} d'où le résultat.

\uline{Corollaire} \unsure{Tout} \deleted{\unsure{schéma}} module quadratique $\Delta$ sur $S$ est canoniquement muni d'une opération de $\mu_2\;(\simeq(\pm1)_S)$, grâce à la formule précédente. L'opération $x\mapsto-x$ \deleted{\unsure{de}} sur $\Delta$ n'est autre que l'involution canonique de $\Delta$ du § 7.

Décrivons \deleted{\unsure{de suite|explicitement}} \inserted{\unsure{enfin,|encore,}} dans le cas général, $B$ comme quotient de
\[
\underline{\mathrm{Aff}}(F)\simeq\underline{\mathrm{Aff}}\,\check{\mathbb V}(L)\simeq\underline{\mathrm{Sym}}(L^\vee)\simeq\coprod_{n\geq0} L^{\otimes(-n)}
\]
par l'idéal engendré par le sous-module inversible engendré par
\[
(\underset{\substack{\uparrow\\\Gamma(L^2\oplus L\oplus\mathcal O_S)}}{\underline{\unsure{c-b+1|c+b+1}}})\otimes L^{-2}\ \subset\ \underline O_S+L^{-1}+L^{\otimes(-2)}\tag*{$(93)$}
\]
% slanted
\margin{\unsure{remonter|renvois|remarque} au § 2}
\deleted{alors l'involution $\sigma$ \unsure{correspon|correspondante}} ce qui donne ici
\[
B\simeq\underline{\mathrm{Sym}}(L^\vee)/(1-4D)L^{\otimes-2},
\]
et notons que l'involution canonique est définie par le passage (\deleted{parmi} les sections de $F\subset\check{\mathbb V}(L)$ satisfaisant $\alpha$), par l'équation
\[
\sigma x=b-x\tag*{$(94)$}
\]
(\unsure{remonte à|variante :|remarque :} § 7) ce qui donne ici $x\mapsto-x$, dans le cas $b=0$, qui est celui du th.
% archive sheet number: 49
